\clearpage
\begingroup
\HistoricalHead{235}{R}{WEITERES ZUR BAHNBESTIMMUNG.}
\HistoricalFoot{}{30\textsuperscript{*}}
\SourceLeaf{240}
\begin{center}
[2.]
\end{center}
\noindent{\small [S. 9]}
\begin{center}
\textls[100]{Vorschriften,}\\
um aus\\
dreien Längen in der Bahn $v,v',v''$ und\\
den correspondirenden Abständen $r,r',r''$\\
abzuleiten:\\
die Bestimmungsstücke des Kegelschnitts, nemlich\\
\begin{tabular}{ll}
h[alben] Parameter $p$,&Eccentricität $e$,\\
halbe grosse Axe $a$,&Länge des Aphelium $\omega$,\\
h[albe] kleine Axe $b$.&
\end{tabular}
\end{center}
{\small
\[
\operatorname{tg}\omega=
\frac{\dfrac{\sin\frac12(v''-v')\sin\frac12(v''+v')}{r}
+\dfrac{\sin\frac12(v'-v)\sin\frac12(v'+v)}{r''}
-\dfrac{\sin\frac12(v''-v)\sin\frac12(v''+v)}{r'}}
{\dfrac{\sin\frac12(v''-v')\cos\frac12(v''+v')}{r}
+\dfrac{\sin\frac12(v'-v)\cos\frac12(v'+v)}{r''}
-\dfrac{\sin\frac12(v''-v)\cos\frac12(v''+v)}{r'}}.
\]}
\noindent Der Zähler dieses Ausdrucks wird
\[
=2\frac ep\sin\tfrac12(v''-v')\sin\tfrac12(v'-v)\sin\tfrac12(v''-v)\mathbin{.}\sin\omega,
\]
\noindent der Nenner
\[
=2\frac ep\sin\tfrac12(v''-v')\sin\tfrac12(v'-v)\sin\tfrac12(v''-v)\mathbin{.}\cos\omega.
\]

Auch leitet man, durch eine leichte Betrachtung, hieraus folgende zur Rechnung bequemere Formeln ab
\[
2\frac{er'}p\sin\{\omega-\tfrac12(v+v'')\}\sin\tfrac12(v''-v)
=\frac{r'}r-\frac{r'}{r''}=A,
\]
\[
\begin{aligned}
&2\frac{er'}p\cos\{\omega-\tfrac12(v+v'')\}\sin\tfrac12(v''-v)\\
&\qquad=\left(\frac{r'}r-1\right)\operatorname{cotg}\tfrac12(v'-v)
-\left(1-\frac{r'}{r''}\right)\operatorname{cotg}\tfrac12(v''-v')=B,
\end{aligned}
\]
\noindent woraus man also zugleich $\omega$ und $\frac ep$ ableitet. Man kann aus dieser Formel auch zwei andre ähnliche ableiten, indem man $v'$ und $r'$ mit $v''$ und $r''$, oder mit $v$ und $r$ vertauscht.
\clearpage
\HistoricalHead{236}{L}{THEORETISCHE ASTRONOMIE. NACHLASS.}
\SourceLeaf{241}
\noindent{\small [S. 10]}\par\noindent Endlich findet man $p$ entweder durch eine von folgenden Formeln
\[
\frac1p=\frac1r+\frac ep\cos(v-\omega)
=\frac1{r'}+\frac ep\cos(v'-\omega)
=\frac1{r''}+\frac ep\cos(v''-\omega)
\]
\noindent oder unmittelbar durch diese
\[
\frac1p=\frac{\frac1r\sin(v''-v')+\frac1{r''}\sin(v'-v)-\frac1{r'}\sin(v''-v)}
{4\sin\frac12(v'-v)\sin\frac12\sin(v''-v')\sin\frac12(v''-v)}.
\]
\noindent Setzt man $e=\sin\varphi$, so ist $b=\dfrac p{\cos\varphi}\quad a=\dfrac p{\cos\varphi^2}$\textsuperscript{[*)]}.\begingroup\renewcommand{\thefootnote}{[*)}\footnotetext{In diesen Formeln wurde ein Schreibfehler verbessert.]}\endgroup
\begin{center}\textls[100]{Beispiel.}\end{center}
\[
\begin{array}{rl@{\qquad}rl}
v=&69^0\,14'\,22''77&\log\dfrac1r=&9{,}574\,2704\\[8pt]
v'=&73\;48\;36''17&\log\dfrac1{r'}=&9{,}577\,2579\text{\textsuperscript{[**)]}}\\[8pt]
v''=&78\;26\;32''49&\log\dfrac1{r''}=&9{,}580\,8208
\end{array}
\]
\[
\frac{r'}r=0{,}993\,144\,633\text{\textsuperscript{[**)]}}
\qquad
\frac{r'}{r''}=1{,}008\,237\,624\text{\textsuperscript{[**)]}}
\]
\begingroup\renewcommand{\thefootnote}{[**)}\footnotetext{Nach S. 234 ist $\log\frac1{r'}=9{,}577\,2759$; auch die Werte von $\frac{r'}r$ und $\frac{r'}{r''}$ zeigen Unstimmigkeiten.]}\endgroup
\[
\begin{array}{rl}
\log\left(\dfrac{r'}r-1\right)=&7{,}836\,0307\,n\\[3pt]
\log\operatorname{cotg}\tfrac12(v'-v)=&1{,}398\,9691\\\cline{2-2}
&9{,}234\,9998\,n\\[4pt]
&-0{,}171\,79075\\
&-0{,}203\,66677\\\cline{2-2}
B=&+0{,}031\,87602\\[3pt]
\log B=&8{,}503\,4641\\
\log A=&8{,}178\,7753\,n\\\cline{2-2}
{}[\log\operatorname{cotg}(\omega-\tfrac12(v'+v''))=]&9{,}675\,3112\,n\\[8pt]
\omega-\tfrac12(v+v'')=&-25^0\,20'\,13''58\\
{}[\tfrac12(v+v'')=]&73\;50\;27''63\\\cline{2-2}
\omega=&48^0\,30'\,14''05
\end{array}
\]
\[
\begin{array}{rl}
\log\left(1-\dfrac{r'}{r''}\right)=&7{,}915\,8020\,n\\[3pt]
\log\operatorname{cotg}\tfrac12(v''-v')=&1{,}393\,1182\\\cline{2-2}
&9{,}308\,9202\,n\\[8pt]
\log\tfrac12 B=&8{,}202\,4341\\
\log\cos(\omega-\tfrac12(v+v''))=&9{,}956\,0750\\\cline{2-2}
&8{,}246\,3591\\[3pt]
\log\sin\tfrac12(v''-v)=&8{,}904\,2957\\\cline{2-2}
&9{,}342\,0634\\[3pt]
\log r'=&0{,}422\,7421\\\cline{2-2}
\log\dfrac ep=&8{,}919\,3213
\end{array}
\]
\clearpage
\HistoricalHead{237}{R}{WEITERES ZUR BAHNBESTIMMUNG.}
\SourceLeaf{242}
\noindent{\small [S. 11]}
{\footnotesize\setlength{\arraycolsep}{3pt}
\[
\begin{array}{rlll}
\log\dfrac ep=&8{,}919\,3213&8{,}919\,3213&8{,}919\,3213\\[8pt]
\log\cos M=&9{,}970\,9153&9{,}956\,1861&9{,}937\,7996\quad[M=v-\omega]\\[8pt]\cline{2-4}
&8{,}890\,2366&8{,}875\,5074&8{,}857\,1209\\[8pt]
{}\left[\log\dfrac1r=\right]&9{,}574\,2704&9{,}577\,2579&9{,}580\,8208\\[8pt]\cline{2-4}
{}\left[\log\left(r\mathbin{.}\dfrac ep\cos M\right)=\right]&9{,}315\,9662&9{,}298\,2495&9{,}276\,3001\\[8pt]
{}[\text{Numerus}=]&0{,}206\,9980&0{,}198\,7236&0{,}188\,9297\\[8pt]
{}\left[\log\dfrac rp=\right]&0{,}081\,7066&0{,}078\,7190&0{,}075\,1562\\[8pt]\cline{2-4}
\log\dfrac1p=&9{,}655\,9770&9{,}655\,9769&9{,}655\,9770\\[8pt]
\log\dfrac ep=&8{,}919\,3213&&\\[8pt]\cline{2-2}
\log e=&9{,}263\,3443&\varphi=10^0\,33'\,59''43&\log\cos\varphi=9{,}992\,5724
\end{array}
\]}
\[
\begin{array}{rl}
p\;\text{. . . . .}&0{,}344\,0230\\
b\;\text{. . . . .}&0{,}351\,4506\\
a\;\text{. . . . .}&0{,}358\,8782.
\end{array}
\]
\clearpage
\endgroup
